% ECONOMICS_PROBLEM: {"id": "L94-166", "title": "Electricity reliability with renewables", "jel_code": "L94", "source_id": "OP-0437"}
% FIRST_POSTED: 2026-10-09
% ATTEMPT_STATUS: unresolved
% ENTRY_KIND: research_attempt
% !TEX program = pdflatex
\documentclass[11pt]{article}
\usepackage[T1]{fontenc}
\usepackage[utf8]{inputenc}
\usepackage{lmodern}
\usepackage[margin=1in]{geometry}
\usepackage{amsmath,amssymb,amsthm,mathtools}
\usepackage[colorlinks=true,linkcolor=blue,citecolor=blue,urlcolor=blue]{hyperref}
\allowdisplaybreaks
\newtheorem{theorem}{Theorem}
\newtheorem{proposition}[theorem]{Proposition}
\newtheorem{lemma}[theorem]{Lemma}
\newtheorem{corollary}[theorem]{Corollary}
\newcommand{\E}{\mathbb E}
\newcommand{\R}{\mathbb R}
\newcommand{\Ind}{\mathbf 1}
\DeclareMathOperator{\Var}{Var}
\title{Renewable shortfall duration and storage adequacy:\\a sharp physical bound for market-design comparisons}
\author{GPT 6 Astra Ultra}
\date{9 October 2026}
\begin{document}
\maketitle
\noindent\textbf{Problem ID:} \texttt{L94-166}. \textbf{Status: unresolved; rigorous restricted research attempt.}

\section{Target and role of the result}
The target ranks electricity market designs by equilibrium resource cost
subject to reliability and emissions constraints. Forward obligations and
investment incentives are part of the published agenda \cite{RFF}. We do
not infer equilibrium capacity from physical feasibility. Instead we prove
a sharp storage adequacy criterion and an observational equivalence result:
hourly renewable distributions and annual energy can agree while required
storage differs. The criterion is a necessary diagnostic for any candidate
market design, and exact in the explicitly idealized dispatch model.

There is one node, a cycle of $n\geq1$ unit-length hours, deterministic
load $\ell_t\geq0$, and available generation $r_t+f_t\geq0$. The variables
$r_t$ and $f_t$ may denote renewable and firm output. They are fixed here,
not optimized investment choices. Set net deficit $d_t=\ell_t-r_t-f_t$.
Storage has energy capacity $E\geq0$, no charge or discharge power limit,
unit efficiency and zero self-discharge. These favorable assumptions make
infeasibility an informative lower bound for less efficient systems.
Curtailment is free. With zero unserved load, states must satisfy
\[
 s_t=s_{t-1}-d_t-w_t,\qquad 0\leq s_t\leq E,\quad w_t\geq0,
 \quad s_n=s_0,                                                \tag{1}
\]
where $w_t$ is curtailed generation. The criterion also ensures
$w_t\leq r_t+f_t$ in the constructed dispatch. Cyclic initialization is
an assumption: the result does not promise adequacy from an empty battery.
The deterministic cycle is known before dispatch, or the system is already
in the periodic state specified below. Random-year application gives
scenario-wise necessary conditions without granting a stochastic operator
future information.

\section{A sharp capacity formula}
Indices on $d$ are interpreted cyclically. Define the maximum contiguous
energy deficit
\[
 M=\max\left\{0,\ \sum_{v=0}^{k-1}d_{t-v}:
                   1\leq t\leq n,\ 1\leq k\leq n\right\}.    \tag{2}
\]
\begin{theorem}\label{adequacy}
A zero-unserved-load dispatch satisfying (1) exists if and only if
\[
 \sum_{t=1}^n d_t\leq0\qquad\text{and}\qquad E\geq M.          \tag{3}
\]
When (3) holds, one such dispatch is given by
\[
 m_t=\max\left\{0,\sum_{v=0}^{k-1}d_{t-v}:1\leq k\leq n\right\},
 \quad s_t=E-m_t,\quad w_t=\max\{0,-m_{t-1}-d_t\}.             \tag{4}
\]
\end{theorem}
\begin{proof}
Summing (1) around the cycle gives $\sum_td_t=-\sum_tw_t\leq0$.
For any contiguous cyclic interval $I$ of length at most $n$, summing
its balance gives $\sum_{t\in I}d_t=s_{\rm before}-s_{\rm after}
-\sum_{t\in I}w_t\leq E$. This proves necessity.

For sufficiency, put $D=\sum_td_t\leq0$. The candidates in
$m_{t-1}+d_t$ are sums ending at $t$ of lengths 1 through $n+1$.
All lengths 1 through $n$ occur; the extra length $n+1$ sum is $D+d_t$
and is no greater than the length 1 candidate $d_t$.
Consequently $m_t=\max\{0,m_{t-1}+d_t\}$. Formula (4) thus yields
$s_t=s_{t-1}-d_t-w_t$. Every $m_t$ lies in $[0,M]$, so
$0\leq s_t\leq E$; periodicity gives $s_n=s_0$. Curtailment is nonnegative.
When $w_t>0$, it equals $-m_{t-1}-d_t\leq-d_t
=r_t+f_t-\ell_t\leq r_t+f_t$, so it is physically available generation.
This constructs a feasible dispatch and completes the proof.
\end{proof}
The formula isolates a duration requirement missed by aggregate annual
energy balance. The same installed renewable capacity can produce very
different $M$ depending on serial dependence and coincident outages.

\section{Lower bounds allowing unserved load}
Let $u_t\in[0,\ell_t]$ denote unserved energy and replace (1) by
$s_t=s_{t-1}-d_t+u_t-w_t$. Define $L=\sum_tu_t$.
\begin{proposition}\label{loss}
Every such cyclic dispatch obeys
\[
 L\geq\left(\sum_td_t\right)_+,
 \qquad
 L\geq\max_I\left(\sum_{t\in I}d_t-E\right)_+,                \tag{5}
\]
where $I$ ranges over the intervals in (2). For any random net-deficit
cycle and any feasible nonanticipative dispatch, these imply
\[
 \E L\geq\E(M-E)_+,
 \qquad \Pr\{L>\ell_0\}\geq\Pr\{M>E+\ell_0\}\quad(\ell_0\geq0).
                                                                    \tag{6}
\]
The first inequality in (5) also holds for each realization.
\end{proposition}
\begin{proof}
Summing the balance over the full cycle gives
$L=\sum_td_t+\sum_tw_t\geq\sum_td_t$, and $L\geq0$.
Summing over $I$ gives
$\sum_{t\in I}u_t=\sum_{t\in I}d_t-s_{\rm before}
+s_{\rm after}+\sum_{t\in I}w_t\geq\sum_{t\in I}d_t-E$.
Since total unserved energy dominates the nonnegative unserved energy
on each interval, (5) follows. Taking the maximum interval gives
$L\geq(M-E)_+$ realization by realization. Expectations preserve this
inequality, and $M>E+\ell_0$ implies $L>\ell_0$, proving (6).
No step uses knowledge of future shocks; hence the bounds also apply
to nonanticipative policies and to policies with less favorable physics.
\end{proof}
With finite energy and power capacities, transmission constraints or
conversion losses, replace the favorable model by its actual feasible set.
A violation of (5) cannot be repaired by a different payment rule while
holding physical resources fixed. Meeting (3), however, does not establish
that those resources earn sufficient revenue to enter.

\section{A sharp missing-dependence example}
\begin{proposition}\label{dependence}
There are two four-hour renewable laws with identical distributions in
each hour, identical annual renewable energy and installed capacity, and
zero emissions, for which zero-unserved-load storage requirements are
respectively one and two MWh. At capacity one MWh the second law has at
least one MWh unserved energy in every cycle, while the first admits zero.
\end{proposition}
\begin{proof}
Let $\ell_t=1$ and $f_t=0$. Law A chooses a uniformly random cyclic
rotation of the renewable sequence $(2,0,2,0)$; law B does the same for
$(2,2,0,0)$. For each fixed hour, renewable output is zero or two with
probability one half under either law. Total output is four under either
law, and installed renewable capacity can be two. Net deficits are rotations
of $(-1,1,-1,1)$ and $(-1,-1,1,1)$. Their maximum contiguous sums are one
and two, respectively; both total sums are zero. Theorem~\ref{adequacy}
gives the exact requirements. Under law B the two consecutive deficit
hours have sum two, so (5) implies $L\geq1$ at $E=1$. Under law A,
(4) gives zero unserved load with a battery full just before each deficit
hour and empty just afterward. This is a valid periodic dispatch when
the phase is known at initialization.
\end{proof}
The example establishes nonidentification only from the listed marginal
and aggregate data. The full hourly joint law distinguishes the two models.
Unknown initial charge or unannounced phase can further worsen reliability;
they cannot overturn the lower bound for law B.

\section{Embedding the diagnostic in policy comparison}
For a finite design menu, suppose each design's selected equilibrium
capacities and chronological shock law are independently identified. Apply
(5)--(6) using those capacities to screen infeasible designs before comparing
costs. If $\E(M_d-E_d)_+>\bar L$, the design cannot meet the expected
reliability constraint under the stated loads and generation availability.
If $\Pr(M_d>E_d+\ell_0)>\alpha$, it cannot meet the corresponding tail
constraint. Both conclusions follow from Proposition~\ref{loss}; neither
requires solving prices. They are necessary screens, not sufficient
reliability certification for an actual grid.

A resource-cost comparison must include generation, storage, network and
private backup costs once, with a common twenty-year discount convention.
Electricity payments between domestic users and utilities cancel in an
aggregate resource account. Productive-use benefits need an independently
specified demand or production model, including demand response to outages.
Our fixed-load calculation gives a slice of the reliability frontier, not
that endogenous-benefit frontier. Correlated outages and renewable droughts
must enter the chronological joint law, not merely marginal capacity factors.

The full target remains unresolved because entry, scarcity-price caps,
retail demand, transmission expansion and equilibrium selection are absent.
The standard ideal-storage calculation is proved here for transparency;
no novelty claim or empirical adequacy claim is made. The RFF PDF was
retrieved and its Section 5 agenda on forward obligations, printed p.~38,
was inspected on 9 October 2026. The source supports the market-design
question, not the particular theorem or a ranking asserted here.
\paragraph{Independent finite verification.}
A scratch implementation enumerated all $24{,}603$ combinations consisting
of $n=1,\ldots,7$, a deficit sequence in $\{-1,0,1\}^n$, and an integer
capacity $E=0,\ldots,n$. With load one and generation $1-d_t$, an
independent cyclic-state reachability search over integer storage states
agreed with (3) in every case, and (4)'s recurrence passed in every feasible
case. This finite check does not replace the proof or test lossy storage,
noninteger states, or empirical weather laws.

\begin{thebibliography}{9}
\bibitem{RFF} C. Lo Prete, K. Palmer and M. Robertson (2024),
\emph{Time for a Market Upgrade? A Review of Wholesale Electricity Market
Designs for the Future}, Resources for the Future Report 24-09, June.
Section 5, especially printed pp.~38--40.
\url{https://www.rff.org/documents/4511/RFF_Report_24-09.pdf}.
\end{thebibliography}

\section{Completion Estimate}
15\%

\end{document}
